\section{Numerical methods}
\label{sec:methods}

The numerical analysis has three stages. First, the finite-precision
reference configurations are placed on the central-configuration manifold
and the equilibria of the infinitesimal body are located independently
(Section~\ref{subsec:eqsearch}). Second, the admissible central-configuration curves are traced
by pseudo-arclength continuation while the equilibrium search is repeated
at every continuation point and the resulting roots are associated into
branches (Section~\ref{subsec:continuation}). Third, changes in equilibrium multiplicity are
localized and tested against the nondegeneracy conditions of a generic
saddle-node fold (Section~\ref{subsec:foldmethods}).

All tables and figures in Sections~\ref{sec:results}--\ref{sec:hill} are generated from the output of the
same numerical implementation. We distinguish throughout between the
arclength coordinate $\ell$ along a central-configuration continuation
curve and the normalized arc parameter $s\in[0,1]$ obtained by rescaling
$\ell$ over an individual admissible arc.


\subsection{Reference configurations and equilibrium search}
\label{subsec:eqsearch}

The reference configurations are specified only to finite precision and
therefore do not lie exactly on the central-configuration manifold; their
initial constraint residuals satisfy
\[
6\times10^{-3}
\lesssim |\psi|
\lesssim
7\times10^{-2}.
\]
For each configuration we hold $(a,b)$ fixed and refine either $c_1$ or
$c_2$, choosing the coordinate that requires the smaller displacement to
reach
\[
\psi(a,b,c_1,c_2)=0.
\]
The scalar refinement is carried out by a bracketed Brent solve. The
resulting configurations satisfy
\[
|\psi|<10^{-12}.
\]
For comparison, the accompanying reproducibility archive reports both the
adopted single-coordinate displacement and the distance to the constrained
nearest point of $\psi=0$ in the $(c_1,c_2)$ plane when both coordinates are
allowed to vary. Every refined configuration is then checked for membership in the
admissible set \eqref{eq:admissible_set}: all three mass ratios are positive,
the denominators in \eqref{eq:mass_ratios} remain nonzero, and no primary
collision occurs.

At a fixed admissible configuration, the equilibria
\eqref{eq:equilibrium_equations} are located by an independent multi-start
Newton search that uses no previously supplied equilibrium coordinates.
For the reference configurations the initial guesses consist of a
deterministic $81\times81$ grid over
\[
[-5,5]^2
\]
together with 2000 uniformly distributed fixed-seed random starts. From
each initial point we apply a damped Newton iteration to
\[
F(x,y)=
\begin{pmatrix}
\Omega_x\\
\Omega_y
\end{pmatrix}
=0,
\]
using the analytic Hessian $H_\Omega$ of
\eqref{eq:hessian} as the Newton Jacobian. A backtracking line search is
used to enforce decrease of $\|F\|$.

Because the effective potential is singular at the primary locations,
initial points and Newton iterates falling within
\[
r_{\min}=10^{-3}
\]
of any primary are rejected. As an additional guard against retaining
solutions numerically attached to a primary singularity, any converged root
within
\[
10r_{\min}=10^{-2}
\]
of a primary is discarded.

The remaining converged solutions are clustered in the $(x,y)$ plane with
tolerance $10^{-6}$ so that repeated convergence to the same equilibrium is
counted only once. Every retained root is required to satisfy
\begin{equation}
\|\nabla\Omega\|<10^{-12}.
\label{eq:equilibrium_residual}
\end{equation}
The equilibrium count is therefore an output of the multi-start search
rather than an input to it. This procedure provides a numerical search for
all equilibria over the stated set of initial conditions; it is not an
analytic completeness theorem. The additional searches used to verify the
eleven-equilibrium regime are described in Section~\ref{subsubsec:eleven}.

Because the analytic derivatives enter both the equilibrium solver and the
subsequent bifurcation and stability calculations, they are independently
validated against central finite differences (step $h=10^{-6}$) at 25
fixed-seed points drawn uniformly in $[-4,4]^2$, each farther than $0.75$
from every primary so that the finite differences are well conditioned. The
analytic gradient
\eqref{eq:potential_gradient} is compared with finite differences of
$\Omega$, and $H_\Omega$ with finite differences of the analytic gradient.
The maximum discrepancies are of order $10^{-9}$ for the gradient and
$10^{-9}$--$10^{-10}$ for the Hessian; independent finite-difference
evaluations of the mixed derivative $\Omega_{xy}$ also agree at this level.

At each equilibrium we compute the spectrum of the variational matrix
\eqref{eq:variational_matrix}. For numerical classification we define
\[
\sigma=\max_j \operatorname{Re}\lambda_j
\]
and use the threshold
\[
\tau_\sigma=10^{-8}.
\]
An equilibrium is classified as spectrally stable when
$\sigma\leq\tau_\sigma$ and the computed spectrum consists of two
numerically purely imaginary conjugate pairs; otherwise it is classified
as unstable. On spectrally stable equilibria the two positive centre
frequencies are recorded as
\[
\omega_1\geq\omega_2>0.
\]
We also record $\det H_\Omega$, which is used as a degeneracy diagnostic in
the fold calculation below, and the equilibrium Jacobi value
\eqref{eq:critical_jacobi}.


\subsection{Central-configuration continuation and branch tracking}
\label{subsec:continuation}

For the continuation study, $(a,b)$ are held fixed and we write
\[
\mathbf q=(c_1,c_2).
\]
The level set
\[
\psi(\mathbf q)=0
\]
is traced by pseudo-arclength continuation. At an accepted point
$\mathbf q_k$ with unit tangent $\mathbf t_k$, the predictor is
\[
\mathbf q^\ast
=
\mathbf q_k+\Delta\ell\,\mathbf t_k,
\]
and the corrector solves
\begin{equation}
\begin{cases}
\psi(\mathbf q)=0,\\[2mm]
\mathbf t_k^{T}(\mathbf q-\mathbf q_k)-\Delta\ell=0.
\end{cases}
\label{eq:arclength_corrector}
\end{equation}
Newton's method is applied to
\eqref{eq:arclength_corrector}, with
$\nabla\psi$ evaluated by central differences using step
$h=10^{-7}$. The tangent is chosen orthogonal to the constraint gradient,
\[
\mathbf t
\propto
\left(-\psi_{c_2},\,\psi_{c_1}\right),
\]
and its orientation is fixed by requiring
\[
\mathbf t_k\cdot\mathbf t_{k-1}>0.
\]
This allows turning points in either $c_1$ or $c_2$ to be traversed without
changing continuation parameter.

The initial continuation step is
\[
\Delta\ell=0.01.
\]
On a failed correction it is halved down to a minimum of $10^{-7}$; after
successful corrections it is increased by a factor of $1.15$, subject to
the maximum value $0.01$. Accepted continuation points satisfy
\[
|\psi|<10^{-13}.
\]

Admissibility is checked after every successful correction. Each component
is followed in both tangent directions until the admissible set is left.
For the arcs encountered in the present study this loss of admissibility
occurs when one of the positive mass ratios tends to zero. The corresponding
boundary is localized by 40 bisections and the vanishing mass ratio is
identified. Each maximal admissible arc is then reparametrized by normalized
arclength
\[
s\in[0,1].
\]

Continuation of already known equilibria is not sufficient to determine
the equilibrium multiplicity, because a pair created at a saddle-node does
not exist on the opposite side of the fold. We therefore repeat an
independent multi-start equilibrium search at every accepted
central-configuration point. For this stage the search set is enlarged to
an $81\times81$ grid over
\[
[-8,8]^2,
\]
1000 fixed-seed random starts, and a $35\times35$ local grid
\[
[x_i-0.8,x_i+0.8]\times[y_i-0.8,y_i+0.8]
\]
around each primary. The local searches are particularly important when a
mass ratio becomes small, because equilibria associated with that primary
can lie close to the corresponding singularity.

The same Newton solver and singularity guards as in Section~\ref{subsec:eqsearch} are used:
initial points and Newton iterates within $10^{-3}$ of a primary are
rejected, converged roots within $10^{-2}$ are discarded, and every retained
root must satisfy \eqref{eq:equilibrium_residual}. Roots are clustered at
$10^{-6}$ within the equilibrium solver; the continuation driver performs
an additional numerical de-duplication before assembling the root set at
each sample.

For branch-dependent quantities, roots at neighbouring continuation samples
are associated by predictor-based one-to-one matching in the equilibrium
$(x,y)$ plane. Suppose that an active branch has its two most recent
positions $\mathbf z_{k-1}$ and $\mathbf z_k$. Its next position is predicted
by linear extrapolation,
\[
\mathbf z_{\rm pred}
=
2\mathbf z_k-\mathbf z_{k-1},
\]
and its most recent inter-sample displacement is
\[
\delta
=
\|\mathbf z_k-\mathbf z_{k-1}\|_2.
\]
A candidate root $\mathbf z$ is eligible for association when
\begin{equation}
\|\mathbf z-\mathbf z_{\rm pred}\|_2
\leq
\max(0.18,\,3\delta).
\label{eq:branch_match_tolerance}
\end{equation}
For a newly initialized branch with only one previous position, the
predictor is that position and the implementation uses the fallback
$\delta=0.1$. Candidate branch--root pairs are ordered by predictor distance
and assigned greedily subject to a one-to-one correspondence.

Because an accepted association may still be geometrically weak, a separate
reliability criterion is imposed. A match is classified as reliable only
when
\begin{equation}
\|\mathbf z-\mathbf z_{\rm pred}\|_2
\leq
\max(0.25,\,2\delta).
\label{eq:branch_reliability}
\end{equation}
Wider accepted matches, together with rapidly escaping roots near a
mass-vanishing boundary, are flagged unreliable and excluded from
branch-dependent stability, frequency, Jacobi, and periodic-orbit claims.
The equilibrium count itself is determined independently by the multi-start
root search and therefore does not depend on branch association. The same
adaptive branch assignment is used for all branch-dependent diagnostics,
figures, and periodic-orbit parent selections reported below.

The constants in \eqref{eq:branch_match_tolerance}--\eqref{eq:branch_reliability}
are implementation choices, and the branch-exchange identification reported in
Section~\ref{sec:results} does not depend on them. Re-tracking the Case-2 arc
with tighter $(0.12,\,2\delta;\ 0.15,\,1.5\delta)$ and looser
$(0.30,\,4\delta;\ 0.40,\,3\delta)$ variants leaves the branch partition through
the eleven-equilibrium interval, and through the entire creation-to-destruction
window of the pair $B_a$, unchanged. In that window the predictor-to-match
distances stay well within the acceptance tolerance
\eqref{eq:branch_match_tolerance} (ratio $\lesssim0.25$ in the eleven-equilibrium
interval, $\lesssim0.8$ across the full seven-fold sequence), and no match in the
fold sequence is flagged unreliable at the default thresholds; the unreliable
flags arise only near the mass-vanishing arc ends. The branch-exchange claim is
therefore insensitive to the threshold choice. These sensitivity data are
provided in the reproducibility archive.


\subsection{Fold localization and verification}
\label{subsec:foldmethods}

Whenever the independent equilibrium search detects a change in
multiplicity between neighbouring continuation samples, the transition is
first localized by seven bisections in the central-configuration parameter,
with the multi-start equilibrium search repeated at every midpoint. The
equilibria present on the higher-count side but absent on the lower-count
side identify the candidate merging pair.

The midpoint of this pair and the corresponding bracketed primary
configuration are then used to initialize a Newton solve of the augmented
system
\begin{equation}
\mathcal G(x,y,c_1,c_2)
=
\begin{pmatrix}
\Omega_x\\
\Omega_y\\
\det H_\Omega\\
\psi
\end{pmatrix}
=
\mathbf 0.
\label{eq:fold_augmented_system}
\end{equation}
Equation~\eqref{eq:fold_augmented_system} consists of four equations in the
four unknowns $(x,y,c_1,c_2)$, with $(a,b)$ fixed. It is solved by damped
Newton iteration using a finite-difference Jacobian with step $10^{-6}$,
and an augmented-system residual not exceeding $10^{-11}$ is required. The
finite-difference Jacobian here differentiates the coupled residual
$(\Omega_x,\Omega_y,\det H_\Omega,\psi)$ with respect to both the equilibrium
coordinates and the central-configuration variables; this outer solve is
distinct from the periodic-orbit correction of
Section~\ref{subsec:periodic_shooting}, which uses the exact analytic
variational matrix \eqref{eq:variational_matrix} integrated as the
state-transition matrix.
Thus the critical equilibrium is simultaneously placed on the
central-configuration manifold and at a degeneracy of the equilibrium map.

Vanishing of $\det H_\Omega$ alone does not establish a generic
saddle-node. Let $\mu_0$ denote the eigenvalue of $H_\Omega$ having the
smallest absolute value at the critical point, let $\mu_\perp$ denote the
other eigenvalue, and let $\mathbf v$ be a unit vector spanning the null
direction,
\[
H_\Omega\mathbf v=0.
\]
A simple fold requires
\[
\mu_0=0,
\qquad
\mu_\perp\neq0,
\]
together with nonzero transversality and quadratic-nondegeneracy
coefficients
\begin{equation}
\alpha
=
\mathbf v^{T}
\left.
\frac{d}{d\ell}
F\!\left(\mathbf z_c;\mathbf p(\ell)\right)
\right|_{\ell=\ell_c},
\qquad
\beta
=
\frac{1}{2}
D^3\Omega(\mathbf z_c)
[\mathbf v,\mathbf v,\mathbf v],
\label{eq:fold_coefficients}
\end{equation}
where $\mathbf z_c=(x_c,y_c)$ is the critical equilibrium and
$\mathbf p(\ell)$ denotes the primary configuration along the
central-configuration curve. In the parameter derivative defining
$\alpha$, the equilibrium position $\mathbf z_c$ is held fixed while the
primary geometry and masses vary consistently along the CC curve through
\eqref{eq:mass_ratios}. Both $\alpha$ and $\beta$ are evaluated by centred
finite differences.

For
\[
\alpha\neq0,
\qquad
\beta\neq0,
\]
Lyapunov--Schmidt reduction of the equilibrium equations onto the null
direction gives the local fold normal form
\[
\alpha(\ell-\ell_c)+\beta u^2+\cdots=0,
\]
and hence
\[
u_\pm
\sim
\pm
\sqrt{-\frac{\alpha}{\beta}(\ell-\ell_c)}.
\]
The separation $d$ of the merging pair must therefore satisfy
\begin{equation}
d\propto|\ell-\ell_c|^{1/2}.
\label{eq:fold_scaling}
\end{equation}
We test this prediction independently by resolving the merging pair on the
side of the fold on which it exists and fitting
\[
d\propto|\ell-\ell_c|^\gamma
\]
on the innermost resolved points in log--log coordinates. Since $s$ is a
smooth monotone rescaling of $\ell$ on each admissible arc, the exponent is
unchanged when expressed in terms of $|s-s_c|$.

Finally, the variational spectrum is evaluated along every reliably tracked
equilibrium branch. Where two centre frequencies are present, the ratio
$\omega_1/\omega_2$ is compared with the low-order rational ratios examined
in Section~\ref{sec:stability}, using the numerical proximity tolerance $0.04$. These
near-resonance flags are indicators for the subsequent periodic-orbit
analysis; they are not treated as evidence of exact nonlinear resonance.